\subsection{Directed dependency graph: step-by-step workflow}
\label{sec:directed_dependency_stepwise}

The directed dependency graph is treated as a separate branch from the undirected
semantic-similarity graph. The reason is simple: a prerequisite relation must respect
curricular time, whereas the current undirected workflow collapses repeated occurrences of the
same harmonized skill into one node within an analytical unit. For the directed branch, the
primary substantive object is the local trajectory DAG, while higher-level graphs are treated as
aggregate support summaries.

\paragraph{Fixed design rules.}
\begin{itemize}
    \item A node may point only to a node that appears later in the same trajectory.
    \item A skill in year $1$ cannot depend on a skill in year $2$ of the same trajectory.
    \item Same-year, same-class pairs are not directed unless a validated within-year ordering
    variable exists.
    \item Dependencies are learned globally across trajectories, but exported locally as one graph
    per trajectory unit.
\end{itemize}

\begin{enumerate}

    \item \textbf{Step D1: Build the ordered trajectory table.}

    \textit{Input.} The harmonized occurrence-level dataset.

    \textit{Action.} Assign each occurrence to exactly one ordered curricular path
    \texttt{trajectory\_unit}. The default rule is:
    \[
    \texttt{trajectory\_unit}
    =
    \begin{cases}
    \texttt{programme} \times \texttt{section}, & \text{if section defines a distinct path},\\
    \texttt{programme}, & \text{otherwise.}
    \end{cases}
    \]
    Retain, at minimum:
    \texttt{trajectory\_unit}, \texttt{year}, \texttt{harmonized\_skill\_id},
    \texttt{occurrence\_id}, and \texttt{class\_id} if available.

    \textit{Expected outcome.} Every occurrence has one valid trajectory assignment and one valid
    time position.

    \textit{Main outputs.}
    \texttt{10\_directed\_trajectory\_panel.csv},
    \texttt{10\_directed\_trajectory\_audit.csv}.

    \item \textbf{Step D2: Create time-layered nodes.}

    \textit{Input.} The ordered trajectory table from Step D1.

    \textit{Action.} Define the directed-graph node as
    \[
    v = (\texttt{trajectory\_unit}, \texttt{year}, \texttt{harmonized\_skill\_id}).
    \]
    If a validated within-year position variable exists, refine the node definition to
    \[
    v = (\texttt{trajectory\_unit}, \texttt{year}, \texttt{position},
    \texttt{harmonized\_skill\_id}).
    \]
    Collapse repeated occurrences only within the same time slice. Do \emph{not} collapse the same
    skill across different years.

    \textit{Expected outcome.} The same skill may appear as different nodes in different years,
    which preserves the temporal structure required for a prerequisite graph.

    \textit{Main outputs.}
    \texttt{11\_directed\_nodes.csv},
    \texttt{11\_directed\_node\_provenance.csv}.

    \item \textbf{Step D3: Split temporally admissible pairs from same-time pairs.}

    \textit{Input.} The time-layered node table from Step D2.

    \textit{Action.} For each pair of nodes within the same trajectory, classify the relation as:
    \begin{itemize}
        \item earlier-to-later;
        \item same year and same class;
        \item same year and different class;
        \item reverse-time impossible.
    \end{itemize}
    Only earlier-to-later pairs enter the prerequisite candidate set.
    Same-time pairs are routed to a separate co-requisite layer.
    Reverse-time pairs are excluded.
    If a validated within-year order variable exists, same-year pairs may be reclassified into
    earlier-to-later within the year.

    \textit{Expected outcome.} The directed branch contains only temporally admissible candidate
    edges, while same-time relations are preserved without being forced into arbitrary direction.

    \textit{Main outputs.}
    \texttt{12\_prerequisite\_candidates.csv},
    \texttt{12\_corequisite\_candidates.csv}.

    \item \textbf{Step D4: Compute global temporal precedence evidence.}

    \textit{Input.} The ordered trajectory table from Step D1.

    \textit{Action.} For each skill $s$ in trajectory $u$, define its first appearance year:
    \[
    T_u(s) = \min\{y : s \text{ appears in trajectory } u \text{ in year } y\}.
    \]
    For each skill pair $(i,j)$, compute:
    \[
    n^+_{ij} = \#\{u : T_u(i) < T_u(j)\}, \qquad
    n^-_{ij} = \#\{u : T_u(i) > T_u(j)\}, \qquad
    n^0_{ij} = \#\{u : T_u(i) = T_u(j)\}.
    \]
    Also compute the precedence score:
    \[
    \mathrm{Prec}_{ij}
    =
    \frac{n^+_{ij} - n^-_{ij}}{n^+_{ij} + n^-_{ij} + n^0_{ij}}.
    \]

    \textit{Expected outcome.} A global evidence table shows whether a skill tends to appear
    before another skill, after it, or at the same time.

    \textit{Main outputs.}
    \texttt{13\_precedence\_statistics.csv}.

    \item \textbf{Step D5: Build the learning sample for direct parent estimation.}

    \textit{Input.} The trajectory table from Step D1 and the precedence table from Step D4.

    \textit{Action.} For each child skill $j$, each trajectory $u$, and each year $y$, define:
    \[
    Y_{u,y,j}
    =
    \mathbf{1}\{\text{skill } j \text{ first appears in year } y+1\},
    \]
    and for each candidate parent skill $i$:
    \[
    X_{u,y,i}
    =
    \mathbf{1}\{\text{skill } i \text{ has appeared by year } y\}.
    \]
    Restrict candidate parents to skills that are temporally admissible for $j$.
    If semantic information is used in the directed branch, estimate parallel candidate sets under
    each representation
    \[
    r \in \{\texttt{name\_only}, \texttt{concat}\}.
    \]

    \textit{Expected outcome.} Each child skill receives a tractable prediction sample for
    estimating its direct parents rather than all earlier skills indiscriminately.

    \textit{Main outputs.}
    \texttt{14\_dependency\_learning\_sample.csv},
    \texttt{14\_candidate\_parent\_sets.csv}.

    \item \textbf{Step D6: Estimate direct parents with sparse conditional models.}

    \textit{Input.} The learning sample from Step D5.

    \textit{Action.} For each child skill $j$, fit a sparse conditional model predicting first
    appearance from previously observed skills. A practical choice is a sparse logistic model or
    a discrete-time hazard model. Repeat this estimation across bootstrap resamples of trajectory
    units and record the edge-selection frequency:
    \[
    q_{ij}
    =
    \frac{1}{B}
    \sum_{b=1}^{B}
    \mathbf{1}\{i \rightarrow j \text{ selected in bootstrap } b\}.
    \]
    When semantic representations are being compared, compute $q_{ij}^{(r)}$ separately for each
    representation $r$.

    \textit{Expected outcome.} Each admissible edge receives an empirical support probability
    instead of a hand-made yes/no decision.

    \textit{Main outputs.}
    \texttt{15\_direct\_parent\_scores.csv},
    \texttt{15\_bootstrap\_selection\_summary.csv}.

    \item \textbf{Step D7: Select the graph rule automatically.}

    \textit{Input.} The bootstrap edge-support scores from Step D6.

    \textit{Action.} For each candidate edge-support cutoff $\tau$, evaluate held-out prediction of
    first-introduction events. If semantic representations are being compared, select jointly:
    \[
    (r^*, \tau^*)
    =
    \arg\max_{r,\tau}
    \mathrm{CVSkill}(r,\tau).
    \]
    If no representation comparison is used, select only
    \[
    \tau^*
    =
    \arg\max_{\tau}
    \mathrm{CVSkill}(\tau).
    \]
    If no thresholded graph outperforms the null model, do not force a binary graph; retain the
    weighted directed graph with edge weights $q_{ij}$.

    \textit{Expected outcome.} The representation choice and the cutoff choice are made
    automatically from predictive performance rather than from hand-set topology rules.

    \textit{Main outputs.}
    \texttt{16\_directed\_selection\_surface.csv},
    \texttt{16\_directed\_selected\_model.csv}.

    \item \textbf{Step D8: Instantiate one local DAG per trajectory unit.}

    \textit{Input.} The selected model from Step D7 and the time-layered nodes from Step D2.

    \textit{Action.} For each trajectory unit $u$, include a directed edge $i \rightarrow j$ when:
    \begin{enumerate}
        \item the edge is selected by the global rule;
        \item both nodes are present in trajectory $u$;
        \item the temporal order in $u$ is valid.
    \end{enumerate}
    Then apply transitive reduction within each local DAG so that redundant edges are removed.
    Export same-time relations separately as co-requisite edges.

    \textit{Expected outcome.} Each trajectory unit has one interpretable local prerequisite DAG,
    plus an optional same-time co-requisite layer.

    \textit{Main outputs.}
    \texttt{17\_directed\_edges\_trajectories.csv},
    \texttt{17\_corequisite\_edges\_trajectories.csv},
    \texttt{17\_directed\_graph\_summary.csv}.

    \item \textbf{Step D9: Aggregate local DAGs upward for comparison.}

    \textit{Input.} The local trajectory DAGs from Step D8.

    \textit{Action.} Aggregate the local graphs into higher-level directional support graphs for
    sections, tracks, programmes, and pooled. For each aggregated edge $i \rightarrow j$, report
    summary support measures such as:
    \begin{itemize}
        \item frequency across local trajectories;
        \item mean support probability $\bar q_{ij}$;
        \item inclusion rate across local units.
    \end{itemize}
    Local trajectory graphs remain strict DAGs. Aggregated higher-level graphs may contain cycles,
    because they summarize heterogeneous local pathways.

    \textit{Expected outcome.} Higher-level directed graphs become comparable because they all
    inherit the same globally learned dependency rule.

    \textit{Main outputs.}
    \texttt{18\_directed\_edges\_sections.csv},
    \texttt{18\_directed\_edges\_tracks.csv},
    \texttt{18\_directed\_edges\_programmes.csv},
    \texttt{18\_directed\_edges\_pooled.csv},
    \texttt{18\_directed\_aggregation\_summary.csv}.

\end{enumerate}

\paragraph{Minimal interpretation rules.}
\begin{itemize}
    \item The main substantive graph is the local trajectory DAG, not the pooled graph.
    \item The pooled graph is a summary of directional support across local trajectories.
    \item Same-year, same-class pairs belong to the co-requisite layer unless a validated
    within-year order variable exists.
    \item Direction is determined by temporal admissibility and predictive support, not by
    semantic similarity alone.
\end{itemize}